%======================================================================
\section{Remarks and open questions}\label{sec:remarks}
%======================================================================

\subsection{Other archimedean data}

The reduction of Section~\ref{sec:reduction} uses only that the Dirichlet coefficients of $\zeta^2$ are non-negative and supported on the
integers, a set closed under multiplication. The Dedekind zeta functions of $\QQ(\sqrt5)$ and $\QQ(\sqrt{-3})$, with the Gamma factors $\GR^2$ and
$\GC$, share this property, so an analogue of the decoupling holds for their squares; we do not write it out. An analogue of the $\Gamma$-only
integer-critical function for those data would then give an exact magic function, hence slack $\le0$. But Part~I, Proposition~5.11, certifies
strictly positive slack for both fields, at the gap $\xin2$ and at their own gaps ($\kstar\ge1.64\cdot10^{-5}$ and $\kstar\ge1.50\cdot10^{-4}$). So,
granted the decoupling, no such analogue exists for them, and the existence statement of this paper is specific to the data of $\zeta$. In the
construction this specificity is visible in the fit between the Gamma factor $\GR$, the conductor $1$ and the lattice $\ZZ$: the kernel
$K_0(4\pi\sqrt{nx})$ comes from Whittaker functions of eigenvalue $\frac14$ for $\mathrm{SL}_2(\ZZ)$, and $\sin^2(\pi x)$ from a second difference
with step $1$.

\subsection{Consequences for the questions of Part I}

With Theorems~\ref{thm:main-S} and~\ref{thm:main-U}, the uniqueness questions left open in Part~I, \S6.3 (non-integer zero weights, extra
prime-side mass together with extra zeros, infinitely many extra zeros without a summability condition) are settled, since every admissible pair
is $\pz$. Part~I asked whether \condU\ forces an exact magic function without spurious zeros; the function $F$ of Theorem~\ref{thm:main-S} is one,
with $Z(F)=\Zz$. The existence statement \condE\ without RH, the last question of Part~I, is equivalent to RH by Corollary~\ref{cor:main-RH}.
Part~I, Corollary~3.12, characterised \condS\ by atomicity of the admissible zero measures; under RH there is now only one admissible pair.

\subsection{How the function was found}\label{ssec:history}

The input was found by analogy with a holomorphic model case. For a model problem whose kernel is $K_1$ instead of $K_0$,
Viazovska's contour argument with the depth-$2$ quasimodular form $E_4-E_2^2=288\sum_nn\sigma_1(n)q^n$ gives, up to a constant factor, the
function $\sin^2(\pi x)\sum_n\sigma_1(n)\sqrt{nx}\,K_1(4\pi\sqrt{nx})$, which is positive off the integers and vanishes doubly at the positive integers;
it is an object of the type studied by Feigenbaum, Grabner and Hardin, formally the case $d=0$ of their family \cite{FGH21}. Its Eisenstein
analogue, built from the completed Eisenstein series $E^*(\tau,\frac12)$, is $\sin^2(\pi x)\sum_nn\,d(n)K_0(4\pi\sqrt{nx})$: it has the right
kernel, but the constant term of the Eisenstein series makes the corresponding $H$ singular at $t=\pm\frac i2$ (compare the remark after
Lemma~\ref{lem:small-x}). Numerical interpolation at the integer nodes, whose canonical interpolants only semi-converge, and a numerical
construction of $\mathrm{Sym}^2$-equivariant eigenvalue-$\frac14$ functions without constant terms then singled out one candidate, which was
recognised, first numerically and then in closed form, as the spectral derivative of Section~\ref{sec:automorphic}. The proofs were
written for that function, and no step of them uses these computations.

\subsection{What remains open}

\begin{enumerate}[label=(\arabic*),leftmargin=*]
\item \emph{RH}, equivalently \condE, equivalently $\kstar=0$, equivalently $\kOPS=0$ (Corollaries~\ref{cor:main-RH} and~\ref{cor:nogap}).
\item \emph{Uniqueness of the integer-critical function.} We do not know whether $H$ is the only even function in $\Wc_\delta$, with $H(0)=1$, whose
$\Gamma$-only transform vanishes at all integers $\ge2$ and is non-negative beyond the gap.
\item \emph{The prime-power family of Part~I.} Whether the exact Hermite family $\Xi^2P_J(t^2)$ of Part~I converges, and to what, is open
(Part~I, Conjecture~6.1). The transform $\hF$ of our function vanishes doubly at $\xin n$ for all integers $n\ge2$, not only at the prime powers, and $F$ is not obtained
as a limit.
Canonical polynomial interpolants at the integer nodes approach it numerically and then diverge.
\item \emph{A proof of positivity without computer.} In our proof only $H(0)>0$ needs no computation (Proposition~\ref{prop:posdef}). For
$t\ne0$ positivity rests on two certified computations on the trivial-bound box: the moment certificate for $|t|\le10.355$
(Theorem~\ref{thm:moments}) and the numerical inputs of the large-$|t|$ theorem at $T_0=9$ (Proposition~\ref{prop:L-numbers-box}); the original
proofs for the magic functions of sphere packing also rested on certified computations \cite{Via17,CKMRV17,CKMRV22}. Theorem~\ref{thm:an-bounds}
with two one-dimensional integrals, which we evaluated in Arb, reaches $|t|<4.4699$ (Remark~\ref{rem:moments}). A proof without computer would
need hand bounds for more moments (six moments, with their certified values, reach only $|t|\le8.47$) and a hand treatment of the large-$|t|$
inputs, which are integrals of $\Bx$ near its zero at $\sqrt3/2$. Positivity does not depend on the arithmetic of the Kloosterman terms beyond
the trivial bound (Remark~\ref{rem:negctl}).
\item \emph{Other Gamma factors and lattices.} Which archimedean data admit an exact magic function, and through which automorphic objects?
\end{enumerate}
